本地文件解决单节点恢复，复制把同一Redis数据历史送到其他节点，为读取和故障接替提供副本。Redis主节点异步向副本传播修改，副本落后时可以继续追赶。主节点已执行并回复一条写入，副本尚未取得这一位置，是协议中正常可达的窗口。准确说明复制进度，先给一条历史命名，再给这条历史中的位置编号~\cite{redis8102-replication}。

\subsection{两套replid与复制位置}

复制历史使用\code{replid}作标识，\code{offset}标记相应复制字节流的位置。A和B拥有同一历史标识，A在1000、B在1023，表示A可以通过取得后续流追到B。两个节点各自都在1023，却属于不同历史，数字相等不足以证明状态相同。这里的offset按复制流计量，业务购买和业务提交有自己的编号。

主节点用复制backlog保留一段最近流。副本重连时，用\code{PSYNC}提交原历史标识与接续位置。主节点先判断历史是否匹配，再判断请求位置是否在保留范围内。两项都满足，就返回继续复制并补发缺少的后缀。断线太久、写流太快使backlog覆盖了旧位置，或者请求历史已经无法识别，就进入全量同步。

\CEighteenFigure{ch18-11}{PSYNC同时检查历史标识与backlog位置}{fig:18-psync}

副本提升为新主时，生成新的当前replid，同时把旧历史标识及其有效截止位置保存在第二套标识中。其他副本如果还持有旧标识，只要位置落在允许范围且backlog可补齐，就有机会接上新主。截止位置很关键：原主在分区期间可能继续写入自己的后缀，那个后缀并未因为旧replid相同就进入新主历史。

例如，旧历史为H1，R提升时已处理到5000。R生成H2，把H1保留为第二历史，并按协议记录旧历史可接续的截止位置。另一个副本V持有H1且已处理到4800，请求从4801继续，R若保留这段backlog，就能补给V共同前缀后面的流。原主M在隔离期间沿H1写到5200，R却没有5001至5200。第二历史的截止检查排除这段分叉被当作共同历史接续~\cite{redis8102-source}。

backlog大小把复制流量与断连时间联系起来。假设保留1 MiB，断连期间持续产生200 KiB/s的复制流，约五秒就能覆盖这一容量，命令开销和瞬时流量还会改变实际窗口。若期望副本在更长的网络抖动后走部分同步，就按实际复制字节率和余量设置保留范围，或者接受全量同步成本。这里的数值是容量演算，业务请求数要经过命令和payload大小才能换算成字节流量。

沿假设的Redis权威路径，普通断连后，副本追到原C1及结果，调用方重试原S便可读取同一绑定。若切换选中的副本没有收到C1，主角色只拥有它现有的数据历史：同S在这里查空，却可能已经用原平台键成功收款。选择新的主节点以前，必须把候选者的数据前缀与系统已经回复的受理放在同一条故障历史中比较。

\subsection{全量同步与增量衔接}

全量同步先建立共同快照，再接上生成快照以后发生的写流。采用磁盘路径时，主节点后台生成RDB，父进程随后把文件传给副本。快照生成期间的修改继续被保留，副本加载快照后应用增量，最终追上当前位置。只有RDB而没有后续写流时，目标停留在旧时点。

无盘同步改变主节点输出快照的路径：子进程直接把RDB流写向副本连接，省去先写主节点磁盘再传出的步骤。一次传输开始后，新到达的副本可能进入下一轮，因此配置中的等待窗口可以让多个副本共同参与同一轮快照。快照、增量和接管次序保留，数据经过的资源发生变化。

三个速度共同影响恢复时间：主节点继续产生修改的速度、快照传输速度、目标加载与应用增量的速度。快照发送越久，后面需要追赶的增量可能越多；目标一边追赶一边接收新流，应用速度不足时会长期落后。把读流量切到副本前，根据接口允许的陈旧程度检查加载完成及当前位置。

\CEighteenFigure{ch18-12}{全量复制的快照流与同步期间增量}{fig:18-full-sync}

副本如何加载是另一组选项。\code{repl-diskless-load disabled}先把收到的RDB保存到文件再加载；\code{swapdb}在加载新数据期间保留旧数据，成功后交换，代价是同时容纳两份数据及相关开销。当前版本还有\code{flushdb}，在无盘加载前清空原数据，加载失败时原数据已经失去；\code{on-empty-db}只在当前数据为空时选择无盘加载，避免保留两份非空数据~\cite{redis8102-source}。

\CEighteenFigure{ch18-24-diskless-load}{四种加载策略分别决定旧数据何时保留、清空或交换}{fig:18-diskless-load}

无盘发送与无盘加载分别由两端配置，副本是否已有可靠RDB文件也据实际路径判断。\code{swapdb}期间是否继续服务旧数据，受能否识别为同一复制历史等条件影响。承担故障接替的副本把加载状态、复制位置、内存余量和持久策略一起纳入资格判断。同步期间读到旧键，只能说明旧数据的服务路径可用。

\subsection{WAIT、WAITAOF与准入限制}

客户端在某一连接写入完整受理对象后，调用\code{WAIT 1 5000}，要求等待至少一个副本确认该连接之前写入对应的复制位置。返回值是实际已确认的副本数，超时也返回当前数目。实现为每个客户端保存它的写位置，再用副本ACK位置比较，因此写入与等待要保持连接关联。连接池把WAIT发到另一条连接，等待对象就可能不再是刚才那次写入。

\code{WAITAOF}进一步等待AOF同步位置，可以分别指定本地主节点和副本的数量要求。返回的两个数分别说明本地与副本的相应同步确认，要求本地同步时需要本地启用AOF。WAIT观察复制处理位置，WAITAOF观察相应AOF同步位置。在事务和脚本等禁止阻塞的上下文中，命令直接返回当前数量~\cite{redis8102-wait}\footnote{WAIT与WAITAOF的等待对象是同一连接此前写入所对应的位置。副本数量、超时及AOF启用条件由各命令分别检查；返回数量不足时，先前写入可能已经执行。}。

\CEighteenFigure{ch18-13}{同一连接的写位置与两类等待确认}{fig:18-wait}

设主节点已经执行对完整受理对象的\code{SET}，随后WAIT超时并返回0。此时GET可能取得刚写入的记录，因为超时发生在确认等待阶段。调用方保存原S，查询或重试原操作；若改用K2发起新购买，服务器会按新身份受理另一项业务。

\code{min-replicas-to-write}和\code{min-replicas-\allowbreak max-lag}约束主节点在观察到足够数量、ACK足够及时的副本时才接受相应写入。它们把复制健康情况放在准入之前，每条写入位置的确认则由等待命令另外观察。副本刚刚断开时，主节点需要观察连接或ACK变化，准入状态才会更新。该机制缩小某些孤立主节点持续写入的窗口，具体窗口受检测与配置影响。

即使WAIT返回了要求的副本数，故障切换也可能提升另一个可用但落后的副本，或者已确认副本与主节点一同丢失。Redis官方把WAIT定义为改善数据安全的手段，其保证范围由异步复制和实际接替条件共同决定。要求保留每份已回复受理的业务，还要让接管过程选择包含这些受理的可恢复前缀；WAIT返回的数量没有单独完成这一选择。下一段按Sentinel的实际候选筛选说明副本数据怎样进入接替判断~\cite{redis8102-wait}。

\subsection{Sentinel的检测、选举与接替}

Sentinel在普通主副本部署旁维护监控与配置。单个Sentinel根据超时形成主观下线判断SDOWN，再收集其他Sentinel的判断。达到该监控项配置的quorum后，可以形成客观下线ODOWN。quorum控制故障判断需要多少观察者支持，执行故障转移的领导者授权还要求已知Sentinel中的绝对多数，并满足相应quorum条件。三个Sentinel、quorum为1时，正常自动故障转移也要取得已知Sentinel的多数授权~\cite{redis8102-sentinel}。

选出的故障转移领导者筛选可提升副本。连接状态、最近通信和过长断连等因素先决定候选资格，\code{replica-priority 0}排除自动提升。合格者按较低优先级数值、更大复制offset、最后按run ID字典序等规则选择。优先级是显式管理策略，因此offset更大的副本可能排在较高优先级候选之后。

设R1的优先级为100、offset为9000，R2优先级为10、offset为8800，两者都通过连接和新鲜度筛选。排序先看优先级，R2可以被选择。把二者都设为100后，offset较大的R1排在前面。R3即使offset为9100，优先级为0也会被排除。运维据此控制哪些位置允许接替，并核对相应恢复点。

五个已知Sentinel、监控quorum为2时，两个相互可达的观察者可能形成足够的下线报告，但故障转移领导者正常授权至少需要三个Sentinel。把quorum设为4，则四个观察支持及相应授权条件共同约束接替。部署要让预期故障之后的存活侧保有这些观察者。

\CEighteenFigure{ch18-14}{Sentinel从观察故障到发布新主配置}{fig:18-sentinel}

领导者指示目标副本提升，观察其主角色，再安排其他副本跟随，并传播新的主地址与配置。客户端查询Sentinel、更新连接目标，已有旧主连接则继续由实际网络和服务器状态决定。旧主恢复连通后，由Sentinel重新配置成副本，它原有的分叉数据随复制关系被新主历史替换。

考虑旧主M与Sentinel多数及副本R断开，但某个客户端能连接M。R在多数侧提升之前，M可能接受并回复写入w。M的本地AOF甚至已经同步w，R却没有收到。网络恢复后，M降级并跟随R，w可能从最终主历史消失。本地磁盘同步与最终接管历史在这里分离。客户端路由刷新、min-replicas限制和故障域部署都影响窗口，外部付款已经发出的部分则继续以原平台键查询和对账。

\CEighteenFigure{ch18-25-split-brain}{旧主本地保存过的写入，可能没有进入最终接管历史}{fig:18-sentinel-fork}

监控配置也需要持久保存并允许Sentinel更新，发现到的其他Sentinel、当前主地址与配置代次关系到后续多数判断和重启恢复。部署三个进程但把它们放在同一台主机，只增加进程层的观察者数量。主机和磁盘共同失效时，它们属于同一个故障域，第21章的地域部署会继续落实这一条件。
\FloatBarrier
